\section{Formal Security Claims}

\textbf{Claim 1: Hash-chain integrity.}
Assuming SHA-256 is collision resistant, an adversary cannot replace any finalized evidence artifact or ledger metadata entry without changing the corresponding ledger commitment and all downstream commitments, except with negligible probability.

\textit{Argument.}
Each artifact $e_i$ is mapped to $h_i = H(e_i)$, and each ledger entry computes $c_i = H(m_i \parallel h_i \parallel p_i)$, where $m_i$ is metadata and $p_i$ is the previous ledger link. Any undetected replacement would require either a collision in $H(e_i)$ or a collision in the entry commitment. Therefore, tampering propagates to all downstream commitments.

\textbf{Claim 2: Context-bound replay resistance.}
If the registry rejects duplicate commitments and the commitment binds the ledger root, artifact hash, IPFS CID, context hash, and action class, replaying an old anchor cannot create a valid new decision for a different evidence context.

\textit{Argument.}
The commitment is not only a raw hash but a context-bound digest. A replayed commitment either duplicates an existing registry entry and is rejected, or it refers to the same original context. It cannot silently authorize a different action or evidence bundle without changing the commitment.

\textbf{Claim 3: Submitter authorization for AuthV2.2.}
In the revised AuthV2.2 attestor, \texttt{submitDecision} succeeds only for accounts holding \texttt{SUBMITTER\_ROLE}. Unauthorized accounts are rejected before proof validation.

\textit{Argument.}
The publication-target attestor inherits OpenZeppelin \texttt{AccessControl}. The function \texttt{submitDecision} is guarded by \texttt{onlyRole(SUBMITTER\_ROLE)}. The targeted negative test verifies that a non-authorized signer is rejected, while positive tests verify deployer role assignment and administrator role granting.
